\paragraph{Archives.}
The OpenAI scored evaluation pool \cite{lightman} supplies a process-supervised reward model (A) and an outcome-supervised model (B) on the same generations. We follow the original evaluator's attempt denominator of 1,860 and exclude null answers from selection. The second release, AceMath-RewardBench \cite{ace}, supplies paired 72B (A) and 7B (B) reward scores on 64 solutions per problem, eight from each of eight generators. We join the two files by problem index within each subject and require identical ordered answer and grade vectors. The original AceMath evaluator uses 100 random ordered subsets and selects the first maximum, which implements uniform tie breaking in expectation. At $N=8$ our exact subject accuracies for the two models are all within 0.10 percentage points of the corresponding published cells. The seven subjects are analyzed separately. MATH-500 is shared in subject matter and benchmark identity across the two releases. The releases provide different generations and reward models.

\paragraph{Replay.}
We specified $N=8,32,M$ before computing outcomes. Both graders see the same eligible pool. We use the archived binary answer grade throughout. The primary unit combines equal nonempty extracted strings within a problem. Empty or unavailable strings among eligible responses receive separate units. We count mixed-grade groups. Where a mixed group exists, the first-occurrence grade defines a string-level sensitivity analysis and response-level results identify the original graded contrast. No algebraic simplification is applied. We retain all problems, including those with no leverage for a particular budget.

We hide the labels for the purposes of selecting a review order and reveal them to evaluate the stopping rule. The proposed order sorts by $|a_i-b_i|$. A pooled baseline sorts by $\max(a_i,b_i)$, so it reviews likely winners from either verifier. Both use stable input order for priority ties. We also compute the expected interval width for a uniform random permutation at a fixed budget: after $b$ of $m$ units it equals $(1-b/m)\sum_i |d_i|$. All review counts include only grades revealed, assuming unit cost for each exact answer or response. They describe an archive replay. The grader was not rerun and no human time was measured. As a design sensitivity we give reliable adjudication a price of two or five ordinary judgments and allow up to five errors among ordinary judgments.

The study protocol records the initial design, the two-tier cost sensitivity specified before examining outcomes, and two follow-ups prompted by the first replay. The latter address tie priorities and whole-question review costs. Exhaustive tests enumerate candidate subsets, admissible grades, corruptions and small review allocations. The replication package contains all budget rows, including inconclusive error-tolerant comparisons.
